%\documentclass[dvipdfmx,a4paper,12pt]{amsart}
\documentclass[dvipdfmx,a4paper,12pt]{article} % titleとe-mailをコメントアウトする．

%%% Packages %%%
\setlength{\lineskip}{0pt}

% --- 和文フォント設定 (ゴシックを使うなら) --- amsartを使う時はコメントアウト
\renewcommand{\kanjifamilydefault}{\gtdefault}
\usepackage{otf}          % min10を避けるため
\usepackage{pxrubrica}    % 和文ルビ


% --- 基本パッケージ ---
\usepackage{graphicx}
\usepackage[all]{xy}
\usepackage{wrapfig}
\usepackage{pgfplots}
\usepackage{color}
\usepackage[dvipsnames]{xcolor}
\usepackage{latexsym}
\usepackage{setspace}
\usepackage{multirow}
\usepackage{enumerate}
\usepackage{enumitem}
\usepackage{comment}
\usepackage[normalem]{ulem} % \emph の下線化を抑止（cancelと共存）
\usepackage{url}

% --- 数学関連 ---
\usepackage{amsmath,amssymb,amsthm,amsfonts,mathtools}
\usepackage{amscd,dsfont,bigdelim,braket,physics,mathrsfs,bm}

% --- 新規追加 ---
\usepackage{tcolorbox}
\tcbuselibrary{breakable, skins, theorems}

% --- showkeys（常に表示） ---
%\usepackage{showkeys}
%\renewcommand*{\showkeyslabelformat}[1]{%
 % \fbox{\parbox{1.6cm}{\normalfont\tiny\sffamily#1\vspace{6mm}}}%
%}

% --- hyperrefは最後に読み込む ---
\usepackage[dvipdfmx,colorlinks,linkcolor=blue,anchorcolor=blue,citecolor=blue]{hyperref}

%%% レイアウト調整（geometryに統一） %%%
\usepackage[
  top=20mm,        % 上余白
  bottom=30mm,     % 下余白
  left=25mm,       % 左余白
  right=25mm,      % 右余白
  headheight=12pt, % ヘッダー高さ
  headsep=10mm,    % ヘッダーと本文の間
  footskip=32pt,   % 本文とフッターの距離
  includehead,
  includefoot
]{geometry}

%%% 行間調整 %%%
\usepackage{setspace}
\setstretch{1.2}

% --- 段落設定 ---
%\setlength{\parskip}{3pt}
%\setlength{\parindent}{0pt}


%%% 追加（重複なし）パッケージ・設定 %%%

% --- 目次の体裁調整 ---
\usepackage{tocloft}
\renewcommand{\contentsname}{目次} % 日本語化
\setlength{\cftbeforesecskip}{0pt}
\setlength{\cftbeforesubsecskip}{0pt}
\setlength{\cftbeforesubsubsecskip}{0pt}

% --- セクション見出しの体裁調整 ---
%\usepackage{titlesec}
%\titleformat*{\section}{\Large\bfseries}
%\titleformat*{\subsection}{\large\bfseries}
%\titlespacing*{\section}{0pt}{1.5ex plus .2ex minus .2ex}{0.8ex plus .1ex}
%\titlespacing*{\subsection}{0pt}{1.0ex plus .2ex minus .2ex}{0.5ex plus .1ex}

% --- ヘッダー／フッター設定 ---
%\usepackage{fancyhdr}
%\pagestyle{fancy}
%\fancyhf{}
%\rhead{岩井 雅崇}
%\lhead{大阪大学 数学専攻}
%\cfoot{\thepage}



% -- enumerate, itemize行間設定
\setlist[itemize]{itemsep=3pt, parsep=0pt}
\setlist[enumerate]{itemsep=3pt, parsep=0pt}

% --- TikZ 設定 ---
\usepackage{tikz}
\usetikzlibrary{positioning, arrows.meta, fit, calc, backgrounds}
\pgfdeclarelayer{background}
\pgfdeclarelayer{foreground}
\pgfsetlayers{background,main,foreground}

% --- footnote がページをまたがない設定 ---
\interfootnotelinepenalty=10000

% --- 目次に表示する階層の深さ ---
\setcounter{tocdepth}{2}

% --- 日本語目次---
\usepackage{pxjahyper}

%%% Theorem 環境 %%%
\theoremstyle{definition}
\newtheorem{thm}{定理}
\newtheorem{lem}[thm]{補題}
\newtheorem{prop}[thm]{命題}
\newtheorem{cor}[thm]{系}
\newtheorem{claim}[thm]{主張}
\newtheorem{dfn}[thm]{定義}
\newtheorem{rem}[thm]{注意}
\newtheorem{exa}[thm]{例}
\newtheorem{conj}[thm]{予想} % 和文版
\newtheorem{prob}[thm]{問題}
\newtheorem{rema}[thm]{補足}
\newtheorem{dfnthm}[thm]{定義・定理}
\newtheorem{ques}[thm]{問題}
\newtheorem{suma}[thm]{まとめ}

%%%%%%定理などを英語で使うとき%%%%%%%%%%%
\begin{comment}
\newtheorem{thm}{Theorem}[section] 
\newtheorem{theo}[thm]{Theorem}
\newtheorem{cor}[thm]{Corollary}
\newtheorem{prop}[thm]{Proposition}
\newtheorem{conj}[thm]{Conjecture}
\newtheorem*{mainthm}{Theorem}
\newtheorem{deflem}[thm]{Definition-Lemma}
\newtheorem{lem}[thm]{Lemma}
\theoremstyle{definition} 
\newtheorem{defn}[thm]{Definition}
\newtheorem{propdefn}[thm]{Proposition-Definition} 
\newtheorem{lemdefn}[thm]{Lemma-Definition} 
\newtheorem{thmdefn}[thm]{Theorem-Definition} 
\newtheorem{eg}[thm]{Example} 
\newtheorem{ex}[thm]{Example} 
\newtheorem{exa}[thm]{Example} 
\newtheorem{ques}[thm]{Question}
\newtheorem{remin}[thm]{Reminder}
\theoremstyle{remark}
\newtheorem{rem}[thm]{Remark}
\newtheorem{setup}[thm]{Setup}
\newtheorem{obs}[thm]{Observation}
\newtheorem{notation}[thm]{Notation}
\newtheorem{cl}{Claim}
\newtheorem{claim}[thm]{Claim}
\newtheorem{assup}[thm]{Assumption}
\newtheorem{step}{Step}
\newtheorem*{clproof}{Proof of Claim}
\newtheorem{cln}[thm]{Claim}
\newtheorem*{ack}{Acknowledgements} 
\numberwithin{equation}{section}
\newtheorem{case}{Case}
\end{comment}
%%%%%%%%%%%%%%%%%

%%% 新規追加コマンド（逆三角関数など） %%%
\newcommand{\Sin}{\text{Sin}^{-1}} 
\newcommand{\Cos}{\text{Cos}^{-1}} 
\newcommand{\Tan}{\text{Tan}^{-1}} 
\newcommand{\invsin}{\text{Sin}^{-1}} 
\newcommand{\invcos}{\text{Cos}^{-1}} 
\newcommand{\invtan}{\text{Tan}^{-1}} 
\newcommand{\Area}{S}
\newcommand{\vol}{\text{Vol}}
\newcommand{\maru}[1]{\raise0.2ex\hbox{\textcircled{\tiny{#1}}}}
\newcommand{\sgn}{{\rm sgn}}

%%% 既存のコマンド群（統合済み） %%%
\newcommand{\rk}[0]{\operatorname{rk}}
\newcommand{\supp}[0]{\operatorname{Supp}}
\newcommand{\Rad}[0]{\operatorname{Rad}}
\newcommand{\Sha}[0]{\operatorname{Sha}}
\newcommand{\sha}[0]{\operatorname{sha}}
\newcommand{\eend}[0]{\operatorname{End}}
\newcommand{\codim}[0]{\operatorname{codim}}
\newcommand{\nd}[0]{\operatorname{nd}}
\renewcommand{\rank}[0]{\operatorname{rank}}
\newcommand{\degree}[0]{\operatorname{deg}}
\newcommand{\Exc}[0]{\operatorname{Exc}}
\newcommand{\pr}{{\rm pr}}
\newcommand{\id}{{\rm id}}
\newcommand{\Sym}{{\rm Sym}}
\newcommand{\End}[0]{\operatorname{End}}
\newcommand{\Coker}[0]{\operatorname{Coker}}

\newcommand{\Supp}{{\rm Supp}}
\newcommand{\Hom}[0]{\mathscr{H}\!\textit{om}}
\newcommand{\Ext}[0]{\mathscr{E}\!\textit{xt}}
\newcommand{\GL}[0]{\operatorname{GL}}
\newcommand{\SheafHom[1]}{\mathscr{H}\!\!\!\text{\calligra om}_{\,{#1}}}
\newcommand{\PGL}[0]{\mathbb{P}\GL(r,\C)}

\newcommand{\Alb}{{\rm Alb}}
\newcommand{\verti}{{\rm vert}}
\newcommand{\hor}{{\rm hor}}
\newcommand{\univ}{{\rm univ}}
\newcommand{\Tor}{{\rm tor}}
\newcommand{\shaf}{\mathrm{sha}}
\newcommand{\Shaf}{\mathrm{Sha}}
\newcommand{\reg}{{\rm{reg}}}
\newcommand{\sing}{{\rm{sing}}}
\newcommand{\qt}{{\rm{qt}}}
\newcommand\sO{{\mathcal O}}
\newcommand{\Div}[0]{\operatorname{div}}
\newcommand{\ddbar}{dd^c}
\newcommand{\cV}{\mathcal{V}}
\newcommand{\deldel}{\sqrt{-1}\partial \overline{\partial}}
\newcommand{\dbar}{\overline{\partial}}
\newcommand{\I}[1]{\mathcal{I}(#1)}
\newcommand{\Aut}[1]{\mathrm{Aut}(#1)}
\newcommand{\Ker}[1]{\mathrm{Ker}(#1)}
\newcommand{\Image}[1]{\mathrm{Im}(#1)}
\DeclareMathOperator{\Ric}{Ric}
\DeclareMathOperator{\Vol}{Vol}
 \newcommand{\pdrv}[2]{\frac{\partial #1}{\partial #2}}
 \newcommand{\drv}[2]{\frac{d #1}{d#2}}
 \newcommand{\ppdrv}[3]{\frac{\partial #1}{\partial #2 \partial #3}}
\newcommand{\underalign}[2]{\quad \underset{\mathclap{\strut #1}}{#2}\quad}
\newcommand{\polar}{\beta}

\newcommand{\R}{\mathbb{R}}
\newcommand{\Z}{\mathbb{Z}}
\newcommand{\N}{\mathbb{Z}_+} % --- 旧定義を優先
\newcommand{\C}{\mathbb{C}}
\newcommand{\Q}{\mathbb{Q}}
\newcommand{\D}{\mathbb{D}}
\newcommand{\mP}{\mathbb{P}}
\newcommand{\mO}{\mathcal{O}}
\newcommand{\B}{\mathds{B}}
\newcommand{\tl}{\hspace{-0.8ex}<\hspace{-0.8ex}}
\renewcommand{\tr}{\hspace{-0.8ex}>}

\newcommand{\xb}[1]{\textcolor{blue}{#1}}
\newcommand{\xr}[1]{\textcolor{red}{#1}}
\newcommand{\xm}[1]{\textcolor{magenta}{#1}}

\newcommand{\illegible}[1]{\textcolor{red}{[ILLEGIBLE: #1]}}

\renewcommand{\thefootnote}{\arabic{footnote}}

\baselineskip = 15pt
\footskip = 32pt%ここから本文．

\begin{document}
\pagestyle{empty}



\begin{center}
{\Large 演習問題 2025年11月6日(木)} \\

%\vspace{5pt}
%{ \large 提出締め切り 2024年月日(火) 15時10分00秒 (日本標準時刻)}
\end{center}

%\vspace{2pt}
%\begin{flushleft}
%{ \large \underline{学籍番号: \hspace{4cm} 名前  \hspace{8cm} } }
%\end{flushleft}

\begin{center}
 {\large 下の問題を解け. なお解答は配布した解答用紙に解答すること.}
  \end{center}
 ただし解答に関しては答えのみならず, 答えを導出する過程をきちんと記すこと. 
 また解答用紙は1人1枚以上提出すること.
  
  \vspace{11pt}
 
   \begin{enumerate}[label={\large \textbf{第}\arabic*\textbf{問}.}]
\item 行列$A$を次で定めるとき, 以下の問いに答えよ.
 $$
 A = 
 \begin{pmatrix}
 2 &4&-3&8 \\
 3&-1&2&-5 \\
  18&0&2&12
 \end{pmatrix}
 $$
 \begin{enumerate}[label=(\arabic*)]
   \setlength{\parskip}{0cm} % 段落間
  \setlength{\itemsep}{0cm} % 項目間
 \item $A$は何行何列の行列か?
 \item $A$の$(3,2)$成分を答えよ.
  \item $A$の第2行を答えよ.
 \item $A$の第3列を答えよ.
 \end{enumerate}
 
\item  次の行列の計算をせよ.

 \begin{enumerate}[label=(\arabic*)]
   \setlength{\parskip}{0cm} % 段落間
  \setlength{\itemsep}{0cm} % 項目間
\item 
$
 \begin{pmatrix} 5 & -2 \\ 0 & 1 \end{pmatrix}-  3 \begin{pmatrix} 3 & 4 \\ 1 & -2 \end{pmatrix}
$
\item $
\begin{pmatrix} 2 & 0 \\ -1 & 3 \end{pmatrix}\begin{pmatrix} 1 & 2 \\ 0 & -1 \end{pmatrix}
+  \begin{pmatrix} 3 & 1 \\ 2 & 2 \end{pmatrix}\begin{pmatrix} 1 & 0 \\ 2 & -1 \end{pmatrix}
$
\item $
\begin{pmatrix} 5 & 0 \\ 0 & 2 \end{pmatrix}^3 + 3 \begin{pmatrix} 5 & 1 \\ 1 & 2 \end{pmatrix}
$
 \end{enumerate}
 
\item 次の問いに答えよ.
 \begin{enumerate}[label=(\arabic*)]
   \setlength{\parskip}{0cm} % 段落間
  \setlength{\itemsep}{0cm} % 項目間
\item $ \begin{pmatrix} 3 & 1 \\ 2 & 4 \end{pmatrix}$の行列式を求めよ.
\item $ \begin{pmatrix} 2 & 1 \\ -5& 3 \end{pmatrix}$の逆行列を求めよ
\item $ \begin{pmatrix} 100 & 99 \\ 99 & 100\end{pmatrix}$の逆行列を求めよ
 \end{enumerate}
 
\item   
 \begin{enumerate}[label=(\arabic*)]
   \setlength{\parskip}{0cm} % 段落間
  \setlength{\itemsep}{0cm} % 項目間
\item 
$
A=\begin{pmatrix} 4 & 2 \\ 0 & 3 \end{pmatrix}
$とする. $A$を対角化せよ. また$A^n$を$n$を用いて表せ. 
\item $
B=\begin{pmatrix} 13 & -30 \\ 5 & -12 \end{pmatrix}
$とする. $B$を対角化せよ. また$B^n$を$n$を用いて表せ. 
 \end{enumerate}
 
 \begin{flushright}
 {\Large 問題は裏面に続く}
 \end{flushright}

 \newpage
 
\item   $2 \times 2$行列
$
A=\begin{pmatrix}
a& b \\
c& d \\
\end{pmatrix}
$とする. 
 $$
\begin{array}{ccccc}
f_A: &\R^2& \rightarrow & \R^2& \\
&\begin{pmatrix}
x \\ y
 \end{pmatrix} & \longmapsto & 
A
\begin{pmatrix}
x \\ y
 \end{pmatrix}  = 
 \begin{pmatrix}
ax + by \\ cx + dy
 \end{pmatrix}
 &
\end{array}
$$
とおき, $A$を$f_A$に対応する行列という. 次の問いに答えよ. 
 \begin{enumerate}[label=(\arabic*)]
   \setlength{\parskip}{0cm} % 段落間
  \setlength{\itemsep}{0cm} % 項目間
  %\item 45度反時計回りに回転する1次変換$f: \R^2 \to \R^2$に対応する$2 \times 2$行列を求めよ. 
%  \item $x$軸に関しての反転(折り返し)を表す1次変換を$g : \R^2 \to \R^2$に対応する$2 \times 2$行列を求めよ.
  \item 「$x$軸に関しての鏡映(反転)を行い, 135度反時計回りに回転する変換」に対応する$2 \times 2$行列を求めよ.
  \item 「$x$軸に関しての鏡映(反転)を行い, 135度反時計回りに回転し, さらに$x$軸に関しての鏡映(反転)を行う変換」を$g : \R^2 \to \R^2$とする. $g$は「$\alpha$度反時計回りに回転する変換」と同じである. $\alpha$の値を求めよ. 
  ただし$0 \leqq \alpha \leqq 360$とする.
  \end{enumerate}


\item 次の行列の計算を行え.
 
(1)
$
\begin{pmatrix}
 1& 0 & 3\\
 4 &-1&3\\
 0 & 2& 1\\
 \end{pmatrix}
 \begin{pmatrix}
 0 &1&-2 \\
 7&-5&4\\
 1 &0&-2 \\
 \end{pmatrix}
 $
 \quad
(2)
$
\begin{pmatrix}
 2 &-1 \\
 4 &0 \\
 3&-5\\
 \end{pmatrix}
 \left\{
 2
 \begin{pmatrix}
 0 &1&-2 \\
 7&-5&4\\
 \end{pmatrix}
 - 3
  \begin{pmatrix}
 1 &-2&6 \\
 4&-1&5\\
 \end{pmatrix}
\right\}
 $
 
\item  次の行列$A,B,C,D$を次で定義する.
 $$
  A=\begin{pmatrix} %14
 -1 & 2 &-5  \\
 \end{pmatrix} 
 \text{, \,\,} 
B= \begin{pmatrix} %33
 1& 0 & 2\\
 0 & 3 & 0\\
 4 & 0 & 5 \\
 \end{pmatrix} %%32
 \text{, \,\,} 
 C=
  \begin{pmatrix}
 -2 &5 & 3\\
1 &-3&0  \\
 \end{pmatrix}
 \text{, \,\,} 
 D= \begin{pmatrix} %%41
 -4\\
 3 \\
 1
 \end{pmatrix}
 $$
 $AA, AB, AC, AD, BA, BB, BC, BD, CA, CB, CC, CD, DA, DB, DC, DD$の16個の組み合わせのうち, 積が定義されるものを全て求め, その積を計算せよ.
 
\end{enumerate}
 
 \newpage
 
  \begin{center}
 {\Large 解答用紙}
% {\Large 演習問題の解答用紙 2024年1月11日(木) } \\
\end{center}

%\vspace{5pt}
%{ \large 提出締め切り 2024年月日(火) 15時10分00秒 (日本標準時刻)}
%\end{center}

%\vspace{2pt}
\begin{flushleft}
{ \large \underline{学籍番号: \hspace{4cm} 名前  \hspace{9cm}   }  }
\end{flushleft}
 

 \end{document}
